\label{chap.1}
\section{Background and Motivation}
Sixth-generation networks will have to carry very different services over the same radios. Streaming video and extended reality want high throughput (enhanced mobile broadband, eMBB). Factory control and connected vehicles want a delay of a few milliseconds and almost no loss (ultra-reliable low-latency communication, URLLC). Sensor networks want to connect huge numbers of cheap devices that each send very little (massive machine-type communication, mMTC).

Network slicing is the usual answer. The operator cuts the radio and core resources into logical slices and promises each slice a service level agreement (SLA). The difficult part is deciding how big each slice should be, because demand keeps moving and is different in every cell. A fixed split wastes capacity when a slice is quiet and breaks the SLA when it is busy. Learning-based controllers do better, but most of the ones we found in the literature only look at what the network is doing right now, so they react after the demand has already moved.

\section{Problem Statement}
We ask two questions in this thesis.
\begin{enumerate}
	\item Can a slicing controller that is told the demand of the next interval allocate bandwidth better than one that sees only the current demand?
	\item Is a decentralised multi-agent controller a better solution to this task than a single central agent?
\end{enumerate}
We also want to try such controllers in a Digital Twin (DT) before anyone would use them on a real network.

\section{Proposed Approach}
Our pipeline has four stages (Table~\ref{tab:pipeline}). A self-supervised learning (SSL) encoder is trained on unlabelled telemetry and gives a compact description of each observation. A temporal graph neural network (TGNN) predicts the demand of every slice in the next interval; it mixes information from neighbouring cells with GraphSAGE \cite{hamilton2017} and follows each cell through time with a recurrent update. A multi-agent reinforcement learning (MARL) controller, with one agent per cell, then splits the capacity of its cell between eMBB, URLLC, mMTC and an idle share. The Digital Twin is a virtual copy of the network where these allocations are trained, tested and, if possible, approved before they are applied.

\begin{table}[h]
	\centering
	\caption{Stages of the proposed pipeline}
	\label{tab:pipeline}
	\begin{tabular}{|p{2.8cm}|p{6.5cm}|p{4.2cm}|}
		\hline
		\textbf{Stage} & \textbf{Role} & \textbf{Output} \\ \hline
		Self-supervised learning & Learns representations of telemetry from unlabelled data by masking or corrupting features and recovering them & Embedding per node and time step \\ \hline
		Temporal GNN & Forecasts next-interval demand per cell and slice from the cell graph and recent history & Predicted demand \\ \hline
		Multi-agent RL & Each cell agent decides how to split its bandwidth between slices & Slice shares per cell \\ \hline
		Digital Twin & Virtual network that trains, tests and checks allocations and supports what-if scenarios & Evaluation, approval, scenario data \\ \hline
	\end{tabular}
\end{table}

\section{Objectives and Contributions}
\subsection{Objectives}
\begin{enumerate}
	\item Survey the literature on Digital Twins, SSL, graph-based forecasting and MARL for network slicing and identify the research gap.
	\item Collect and prepare public data for each stage of the pipeline.
	\item Implement each module and evaluate it against strong, simple baselines on real data.
	\item Provide one environment in which prediction and allocation are tested together on the same traffic.
	\item Build a working prototype.
\end{enumerate}

\subsection{Contributions}
Each contribution below is supported by experiments in the later chapters, and the limits of each are stated there.
\begin{enumerate}
	\item \textbf{A unified slicing environment.} No public dataset combines topology, per-slice traffic, quality of service (QoS) and allocation actions. We built one from real Telecom Italia Milan traffic, with latency calibrated on a public 6G radio access network (RAN) telemetry dataset, and we mark every component as real, calibrated or modelled.
	\item \textbf{A fair evaluation of MARL.} A parameter-shared controller keeps a fixed-size actor from 9 to 100 cells while a central agent does not scale. Against expert rules tuned on the training days, it wins on all seven test days at every size (by 0.023 to 0.026 in reward), and an on-policy PPO-family learner also beats a rule tuned per time of day and cell class (by 0.009 at 25 cells, on every test day), whereas the off-policy learner only matches that rule.
	\item \textbf{A measurement of what prediction is worth.} Giving the expert rule the TGNN forecast raises its reward by about 0.03 on every test day, about one fifth of the gap to perfect knowledge. The benefit to the learned controller is not established.
	\item \textbf{A city-scale forecasting finding.} A forecaster trained only on the dense centre was 10.8 times worse than persistence in the outer city. Per-cell scaling, random spatial patches and more training days fixed this, and the retrained forecaster beats persistence in all 10{,}000 cells (18\%, 33\% and 31\% lower error for the three slices).
	\item \textbf{A clean negative result for SSL.} After removing duplicate flows and leaked columns, pretraining gave no gain over strong raw-feature baselines at any label budget, on 5G-NIDD and on the Milan forecaster.
	\item \textbf{A Digital Twin that checks decisions.} An approval gate in the environment protects the reward from poor proposals, and a Simu5G load sweep tests the QoS model against a packet-level simulator.
	\item \textbf{A dataset and a prototype.} A 393.6-million-row, 122-column dataset with six what-if scenarios, and the Slice Twin monitor, which is published on the web.
\end{enumerate}

\subsection{Scope}
The thesis does not cover the following.
\begin{itemize}
	\item A packet-level simulator inside the training loop. The QoS comes from a calibrated surrogate, and the telemetry used to calibrate it is itself simulated. Our own 19-gNodeB scenario is not connected to Simu5G; we only ran a single-cell load sweep.
	\item A statistically wide evaluation. Results use one held-out week (seven days), at most three seeds per learned controller (one at 100 cells), and one city.
	\item Deployment on a live or emulated network. The prototype replays recorded days.
\end{itemize}

\subsection{Organisation}
Chapter~\ref{chap.2} reviews the literature and the research gap, and Chapter~\ref{chap.3} describes the datasets and how we prepared them. Chapter~\ref{chap.4} presents the system design, the problem formulation and the unified environment. Chapters~\ref{chap.5}, \ref{chap.6} and \ref{chap.7} cover self-supervised representation learning, temporal GNN forecasting and MARL for slice allocation. Chapter~\ref{chap.8} gives the integrated evaluation, the Digital Twin dataset, the approval gate and the Simu5G check. Chapter~\ref{chap.9} describes the Slice Twin prototype, and Chapter~\ref{chap.10} ends with the limitations and future work.
